
\vspace{-3mm}
\begin{figure*}[ht!]
    \centering
    \resizebox{\textwidth}{!}{
        \begin{forest}
            forked edges,
            for tree={
                grow=east,
                reversed=true,
                anchor=base west,
                parent anchor=east,
                child anchor=west,
                base=center,
                font=\large,
                rectangle,
                draw=hidden-draw,
                rounded corners,
                align=center,
                text centered,
                minimum width=5em,
                edge+={darkgray, line width=1pt},
                s sep=3pt,
                inner xsep=2pt,
                inner ysep=3pt,
                line width=0.8pt,
                ver/.style={rotate=90, child anchor=north, parent anchor=south, anchor=center},
            },
            where level=1{text width=13em,font=\normalsize,}{},
            where level=2{text width=25em,font=\normalsize,}{},
            where level=3{text width=35em,font=\normalsize,}{},
            where level=4{text width=20em,font=\normalsize,}{},
            [\textbf{Adversarial} \textbf{Attacks} \textbf{in LLMs}, for tree={fill=a4},name=adv
                [\textbf{Jailbreak} \S\ref{jailbreak}, for tree={fill=medium-red}
                    [\textbf{Optimization} \S\ref{optimization}
                        [Societal Harm\cite{wu2024llms, pair23, tap23}]
                        [Privacy Violation\cite{wu2024llms, pair23, tap23}]
                        [Disinformation \& Deception\cite{wu2024llms, pair23, tap23}]
                    ]
                    [\textbf{Long Tail Distribution} \S\ref{long}
                        [Rare Prompts\cite{jiang2023promptpacker}]
                        [Out-of-Distribution Exploits  \cite{schulhoff2023hackaprompt}]
                        [Persuasive Manipulation\cite{jiang2023promptpacker}]
                    ]
                ]
                [\textbf{Control Generation} \S\ref{control}, for tree={fill=light-yellow}
                    [\textbf{Direct Attack} \S\ref{direct}
                        [Malicious Prompt Engineering \cite{jiang2023promptpacker}]
                        [Syntax Manipulation \cite{jiang2023promptpacker}]
                        [Prompt Suffix Exploits \cite{schulhoff2023hackaprompt}]
                    ]
                    [\textbf{Indirect Attack} \S\ref{indirect}
                        [Goal Hijacking \cite{chen2024pseudo}]
                        [Prompt Leaking \cite{li2024pleak}]
                        [External Source Injection \cite{greshake2023indirect}]
                    ]
                ]
                [\textbf{Performance Degradation} \S\ref{performance}, for tree={fill=light-blue}
                    [\textbf{Dataset Poisoning} \S\ref{data}
                        [Label Flipping \cite{greshake2023indirect}]
                        [Data Corruption \cite{greshake2023indirect}]
                        [Poisoned Sample Injection \cite{greshake2023indirect}]
                    ]
                    [\textbf{Prompt Injection} \S\ref{prompt}
                        [Wrong Classification \cite{greshake2023indirect}]
                        [Answer Disparity \cite{greshake2023indirect}]
                        [Consistency Violation \cite{greshake2023indirect}]
                    ]
                ]
            ]
        \end{forest}}
    \caption{
    \textbf{Taxonomy of Adversarial Attacks in LLMs.} A structured classification spanning three principal branches—\textbf{Jailbreak}, \textbf{Control Generation}, and \textbf{Performance Degradation}—each reflecting distinct adversarial intents: bypassing alignment, subverting generation control, or degrading functional reliability. Subtypes distinguish \textit{direct vs. indirect} mechanisms and expose \textit{long-tail vulnerabilities}, including rare prompt exploits and semantic hijacks. Anchored in canonical papers, this taxonomy is a conceptual scaffold for reasoning about threat surfaces, model failure modes, and the generality of alignment defenses across adversarial regimes.
    }
    \label{fig:lit_surv}
    %\vspace{-2mm}
\end{figure*}



\vspace{-3mm}
\begin{figure*}[ht!]
    \centering
    \includegraphics[width=\textwidth]{Figures/LLM_Attack_Benchmark_Heatmap_AdjacentBar.pdf}
    \vspace{-8mm}
    \caption{
    \textbf{Benchmarking LLM Vulnerabilities to Jailbreak Attacks.}  
    This heatmap summarizes \textbf{attack success rates} (\textit{higher is worse}) across diverse jailbreak strategies applied to both open and proprietary LLMs. Each row denotes a distinct \textsc{attack category}, targeting prompt alignment, instruction controllability, or generation stability. Key takeaways:  
    \textbf{(i)} \textbf{LLaMA-3} and \textbf{GPT-4} variants show comparatively stronger refusal behavior across adversarial regimes;  
    \textbf{(ii)} \textbf{Vicuna} and \textbf{phi-series} models are especially susceptible to persona-based threats like \textsc{DAN}, \textsc{TAP}, and \textsc{Puzzler};  
    \textbf{(iii)} \textsc{Prompt Extraction} and \textsc{Goal Hijacking} succeed across model families, exposing generalization gaps in safety alignment;  
    \textbf{(iv)} compositional chains like \textsc{BadChain} and continual-learning exploits (\textsc{advVCL}) reveal progressive alignment erosion.  
    The \textit{right-aligned color bar} encodes success rates from 0 (safe) to 100 (compromised), enabling cross-architectural comparison of robustness.
    }
    \label{fig:llm_heatmap_benchmark}
    \vspace{-4mm}
\end{figure*}







\vspace{-2mm}
\caption{
\textbf{Components of the Neural Virulence Index (\textit{nVI}).} 
Each term corresponds to a biologically inspired mechanism governing adversarial semantic takeover in transformer models:
\begin{itemize}[itemsep=0.3em]
  \item \(\boldsymbol{\mathbb{I}_{t^*}}\): \emph{\textbf{Trigger indicator}}. A binary gate activated only by the presence of a rare adversarial token \(t^*\), analogous to viral \textbf{\textit{tropism}}—selective infection of specific tissues or contexts \cite{pomerantz1990tropism}.
  \item \(\boldsymbol{\psi_\ell}\): \emph{\textbf{Layer susceptibility coefficient}}. Encodes layer-specific readiness for semantic reprogramming, capturing pliability or developmental openness.
  \item \(\boldsymbol{\Delta \kappa_\ell, \Delta \mathcal{T}_\ell}\): \emph{\textbf{Curvature and thermodynamic divergence}}. Quantify geometric and energetic deviations from base model states.
  \item \(\boldsymbol{\text{nTDS}_\ell}\): \emph{\textbf{Neural Total Drift Score}}. Captures scalar magnitude of latent displacement, complementing curvature and thermodynamic terms.
  \item \(\mathbf{nDIV}_\ell\): \emph{\textbf{Directional Inheritance Vector}}. Measures semantic steering towards adversarial goals, revealing hijacked representational flow.
  \item \(\mathbf{nCCL}_\ell\): \emph{\textbf{Cultural Conflict Loss}}. Quantifies semantic discord between attacked and base states.
  \item \(\mathbf{nEPI}_\ell\): \emph{\textbf{Epistemic Plasticity Index}}. Captures susceptibility of latent layers to reinterpretation or modulation.
\end{itemize}
Further, the \textit{nVI} can be interpreted as the product of \textbf{thermodynamic drift} and \textbf{semantic virulence}, gated by the trigger token presence:
\[
\boxed{
\text{nVI}(t^*) = \sum_{\ell = \ell_s}^{\ell_e}
\mathbb{I}_{t^*} \cdot 
\left( \text{Drift}_\ell \cdot \text{Hijack}_\ell \right),
}
\]
where
\[
\text{Drift}_\ell := \lambda_\kappa \cdot |\Delta \kappa_\ell| + \lambda_T \cdot |\Delta \mathcal{T}_\ell| + \lambda_{\text{tds}} \cdot \text{nTDS}_\ell,
\quad
\text{Hijack}_\ell := \lambda_{\text{div}} \cdot \text{nDIV}_\ell + \lambda_{\text{conf}} \cdot \text{nCCL}_\ell + \lambda_{\text{epi}} \cdot \text{nEPI}_\ell.
\]
\textbf{Interpretation.}  
This biologically inspired formulation emphasizes that neural \textbf{\emph{semantic infection}} requires two key conditions: (1) a measurable \textbf{energetic/geometric drift} from baseline (Drift), and (2) \textbf{vulnerable semantic channels} (Hijack) amenable to adversarial manipulation. The product ensures that \emph{pathogenic effects manifest only when both access and susceptibility coexist}, echoing classical virulence-host susceptibility models in biology.
}
\label{fig:nVI_equation}
\vspace{-2mm}
\end{figure*}



\begin{itemize}[itemsep=1.2em]

  \item \textbf{Susceptibility (\emph{Semantic Tropism})}:  
  An attack only takes hold if the model enters a \textbf{\emph{receptive state}}—most commonly in mid-depth layers ($\ell \approx 24$–$27$) where \textbf{\textit{epistemic plasticity}} is high. These layers behave like \emph{semantic stem zones}: cognitively pluripotent, weakly canalized, and easily reprogrammed. This mirrors \textbf{\textit{tissue tropism}} in virology, where only certain cell types—those with open chromatin or exposed surface receptors—permit infection \cite{pomerantz1990tropism, zhang2021epigenetic, frantz2015cell}. Without sufficient pliability, even structurally toxic prompts are ignored by the model’s internal logic.

  \item \textbf{Activation (\emph{Latent Regulatory Trigger})}:  
  The adversarial input must align with the model’s internal routing in a way that \textbf{\emph{activates}} dormant behavioral machinery. This is analogous to \textbf{\textit{proviral activation}}, where integrated viral DNA lies silent in the genome until a specific stressor or signaling cascade reawakens it \cite{temin1974provirus, best1996activation, grow2015intrinsic}. In the LLM setting, the trigger token $t^*$ functions as a \emph{semantic ligand}—harmless in isolation, but catalytically potent when presented in the correct context. Recent work confirms this structure: prompt injections only succeed when embedded at \textit{precisely the right semantic junction}, akin to finding an open promoter in chromatin \cite{zou2023universaltransferableadversarialattacks, liu2023jailbroken}.

  \item \textbf{Inheritance (\emph{Downstream Semantic Flow})}:  
  Lasting disruption requires that the adversarial signature be \textbf{\emph{preserved, amplified, and inherited}} across depth. This is quantified by \textbf{nDIV}$_\ell$—the \textit{directional inheritance vector}—which tracks how representational flow is bent from its midpoint. Biologically, this parallels \textbf{\textit{epigenetic memory}}: once an infection alters transcriptional pathways or chromatin marks, the modified state persists across cell divisions \cite{jaenisch2003epigenetic, bird2007perceptions, cedar2009epigenetics}. In transformers, residual and attention mechanisms act as the \emph{semantic cytoskeleton}, enabling adversarial signals to propagate and solidify \cite{hendrycks2021aligning, zhu2024promptbench}.

\end{itemize}

\vspace{0.5em}
\noindent
\textbf{In short:} LLM attacks operate not as brute distortions, but as \textbf{\emph{semantic infections}}—strategically exploiting the model’s internal pliability, latent receptors, and propagation mechanisms. Much like a virus, an adversarial token $t^*$ is only \emph{pathogenic} when three biological-style constraints are met: \textbf{\emph{access}}, \textbf{\emph{activation}}, and \textbf{\emph{inheritance}}. Without all three, the attack fails silently.


